\documentclass[12pt,a4paper]{article}

% Packages
\usepackage[utf8]{inputenc}
\usepackage[margin=0.75in]{geometry}
\usepackage{enumitem}
\usepackage{titlesec}
\usepackage{xcolor}
\usepackage{hyperref}
\usepackage{fontawesome5}
\usepackage{tabularx}
\usepackage{multirow}
\usepackage[T1]{fontenc}
\usepackage{libertine} % Modern, elegant font
\usepackage{xeCJK} % Support for Chinese characters
\setCJKmainfont{KaiTi} % Set Chinese font to KaiTi (楷体)

% Color scheme
\definecolor{primarycolor}{RGB}{72, 61, 139}
\definecolor{accentcolor}{RGB}{70, 130, 180}
\definecolor{textcolor}{RGB}{40,40,40}
\definecolor{linkcolor}{RGB}{70, 130, 180}

% Hyperlink setup
\hypersetup{
    colorlinks=true,
    linkcolor=linkcolor,
    urlcolor=linkcolor,
    citecolor=primarycolor
}

% Section formatting
\titleformat{\section}
  {\color{primarycolor}\Large\bfseries\scshape}
  {}{0em}{}[\color{primarycolor}\titlerule]
\titlespacing{\section}{0pt}{12pt}{6pt}

\titleformat{\subsection}
  {\color{accentcolor}\large\bfseries}
  {}{0em}{}
\titlespacing{\subsection}{0pt}{8pt}{4pt}

% Custom commands
\newcommand{\cvitem}[4]{
  \textbf{#1} \hfill \textit{#2} \\
  \textit{#3} \hfill #4 \\[4pt]
}

\newcommand{\publication}[1]{
  \item #1
  % \textit{\small Contribution: #2}
}

% Remove page numbers
\pagestyle{empty}

% Adjust spacing
\setlength{\parindent}{0pt}
\setlength{\parskip}{0pt}

\begin{document}

% Header
\begin{center}
  {\Large\bfseries\color{primarycolor} Curriculum Vitae - Gaoxiang Jin  金高翔} \\[10pt]
  {\color{textcolor} PhD student \ $|$ \ Max Planck Institute for Astrophysics} \\[10pt]
  
  \begin{tabular}{c c c}
    \faIcon{envelope} \href{mailto:gxjin@mpa-garching.mpg.de}{gxjin@mpa-garching.mpg.de} &
    % \faIcon{phone} +1 (123) 456-7890 &
    % \faIcon{home} Munich, Germany
    \faIcon{globe} \href{https://wwwmpa.mpa-garching.mpg.de/~gxjin/}{wwwmpa.mpa-garching.mpg.de/$\sim$gxjin} &
    % \faIcon{github} \href{https://github.com/yourusername}{github.com/yourusername} &
    \faIcon{orcid} \href{https://orcid.org/0000-0003-3087-318X}{ORCID}
  \end{tabular}
\end{center}

\vspace{1pt}

% Education
\section*{Education}

\cvitem{Ph.D. in Astrophysics}{Expected May 2026}{Max Planck Institute for Astrophysics}{Garching, Germany}
\vspace{-2pt}
\textit{Thesis:} Spatially Resolved Study on AGN Host Galaxies \\
\textit{Advisor:} Prof. Guinevere Kauffmann

\vspace{8pt}

\cvitem{M.S. in Astrophysics}{June 2022}{National Astronomical Observatories of Chinese Academy of Sciences}{Beijing, China}
\vspace{-2pt}
\textit{Thesis:} Active Galactic Nuclei in Galaxy Mergers \\
\textit{Advisor:} Prof. Y. Sophia Dai

\vspace{8pt}

\cvitem{B.S. in Astronomy}{June 2019}{Nanjing University, School of Astronomy and Space Science}{Nanjing, China}
\vspace{-2pt}
% \textit{Advisor:} Prof. Yong Shi


\section*{Research Interests}
Observational astrophysics, primarily in: Co-evolution of AGN and galaxies, Extragalactic gas ecosystems, Baryon cycle in galaxy evolution, Star formation and quenching of galaxies.


\section*{Telescope Time}

\begin{itemize}[leftmargin=*, itemsep=2pt]
  \item \textit{Palomar 200-inch Hale Telescope}, PI, 4 nights, 2022, optical IFU (Cosmic Web Imager)
  \item \textit{Xinglong 2.16 meter Telescope}, PI, 0.5 night, 2020, long-slit spectrograph
  \item \textit{Atacama Large Millimeter/submillimeter Array}, CoI, 12 hours, 2025, ACA
  \item \textit{James Clerk Maxwell Telescope}, CoI, 33 hours, 2022, submillimeter spectrum
  \item \textit{Five-hundred-meter Aperture Spherical radio Telescope}, CoI, 43 hours, 2020/2025, spectrum
  \item \textit{Giant Metrewave Radio Telescope}, CoI, 20 hours, 2022, radio spectrum
\end{itemize}

% Skills
\section*{Technical Skills}

\begin{itemize}[leftmargin=*, itemsep=2pt]
  \item\textbf{Programming:} Python (NumPy, SciPy, Matplotlib, Astropy, scikit-learn), IDL, Shell
  \item \textbf{Data Analysis:} Integral field spectroscopy (optical to infrared spectra fitting; stellar population synthesis, AGN spectra fitting), SED fitting (familiar with CIGALE, Kcorrect, prospector; experience on photometry+spectral fitting), photometry analysis (familiar with SExtractor for optical and IR images, PyBDSF for radio images); building multiwavelength catalog for large surveys (cross-matching archive data); post-processing of simulation data.
  \item \textbf{Observational data reduction:} MaNGA (IFU data), LOFAR (radio interferometry imaging), FAST (21-cm spectra), CWI (optical IFU),
  SDSS (optical photometry and spectra), Subaru-HSC (deep optical imaging), Spitzer (infrared photometry) 
  \item \textbf{Languages:} English (fluent), Mandarin Chinese (native)
\end{itemize}


\section*{Service and Outreach}

\begin{itemize}[leftmargin=*, itemsep=2pt]
  \item \textbf{Conference LOC}, The multiscale environment of AGN across cosmic time (2025), From AGN to starburst -- a multi-wavelength synergy (2019)
  \item \textbf{Journal referee}, The Astrophysical Journal
  % \item \textbf{Mentor of junior students}, Yilu Tong on galaxy merger classifications (2021-2023), and Ludmila Schneider on catalog cross-matching (2022)
  \item \textbf{Co-organizer and Volunteer}, Open Day 2024 at the Garching Research Campus
  \item \textbf{Co-organizer and Volunteer}, Girls Day 2023 at the Max Planck Institute for Astrophysics
\end{itemize}


\section*{Contributed talks}

\begin{itemize}[leftmargin=*, itemsep=4pt, parsep=0pt]

  \item \textbf{The multiscale environment of AGN across cosmic time} \hfill \textit{Munich, Germany, Nov 2025}

  \item \textbf{MIAPbP: The next Generation of AGN Surveys} \hfill \textit{Garching, Germany, March 2025}
  
  \item \textbf{European Astronomical Society Meeting} \hfill \textit{Padova, Italy, July 2024}

  \item \textbf{LOFAR Family Meeting} \hfill \textit{Leiden, Netherlands, June 2024}

  \item \textbf{Resolving Galaxy Ecosystems Across All Scales} \hfill \textit{Hong Kong, China, Dec 2023}
  
  \item \textbf{LOFAR Family Meeting} \hfill \textit{Olsztyn, Poland, June 2023} 

  \item \textbf{The Multiphase Circumgalactic Medium} \hfill \textit{Munich, Germany, March 2023} 

  \item \textbf{SDSS 2021 Collaboration Meeting} \hfill \textit{Online, Aug 2021} 

  \item \textbf{Month of MaNGA 2020} \hfill \textit{Online, Nov 2020} 
  
\end{itemize}

\section*{Summer Schools}
\begin{itemize}[leftmargin=*, itemsep=4pt, parsep=0pt]
  \item \textbf{Fundamentals of Deep Learning} \\ \hfill \textit{Garching, Oct 2025}
  \item \textbf{ESO-Gruber summer school: From Nearby Worlds to Distant Galaxies} \\ \hfill \textit{Garching, June 2025}
  \item \textbf{18th Heidelberg Summer School: Unraveling Galaxy Evolution with JWST} \\ \hfill \textit{Heidelberg, Sep 2023}
  \item \textbf{Frontend research at low radio frequency Radio astronomy} \\ \hfill \textit{online, March 2023}
  \item \textbf{International Summer School on the ISM of Galaxies} \\ \hfill \textit{online, July 2021}
\end{itemize}

\section*{Awards \& Honors}

\begin{itemize}[leftmargin=*, itemsep=2pt]
  \item {\bf National Graduate Scholarship}, University of Chinese Academy of Sciences (2021)
  \item {\bf Outstanding Volunteer Award}, Red Cross Society of China (2017)
  \item {\bf Top-notch Student Scholarship}, Nanjing University (2015)
\end{itemize}


\newpage

% Publications

\section*{Publication list}

\subsection*{First Author Publications}

\begin{enumerate}[leftmargin=*, itemsep=6pt]
  \item \textbf{``The Host Galaxies of Radio AGN: New Views from Combining LoTSS and MaNGA Observations''} \\
  \textbf{Gaoxiang Jin}, Guinevere Kauffmann, Philip N. Best, et al. 2025, \textit{A\&A, 694, A309}
  
  \item \textbf{``Comparison of Global HI and H$\alpha$ Line Profiles in MaNGA Galaxy Pairs with FAST''} \\
  \textbf{Gaoxiang Jin}, Y. Sophia Dai, Cheng Cheng, et al. 2025, \textit{ApJ, 980, 267}

  \item \textbf{``An IFU View of the Active Galactic Nuclei in MaNGA Galaxy Pairs''} \\
  \textbf{Gaoxiang Jin}, Y. Sophia Dai, Hsi-An Pan, et al. 2021, \textit{ApJ, 923, 6}

  \item \textbf{``A Spatially Resolved Evolutionary Sequence of Multi-wavelength Selected AGN Host Galaxies''} \\
  \textbf{Gaoxiang Jin}, Guinevere Kauffmann, Michael R. Blanton, et al. \textit{Submitted to MNRAS}
\end{enumerate}

\subsection*{Co-Authored Publications}

\begin{enumerate}[leftmargin=*, itemsep=6pt]
  \item \textbf{``The origin of double-peaked narrow emission-line galaxies in MaNGA Survey''} \\
  Zhiyun Zhang, [5 authors], {\bf Gaoxiang Jin} et al. 2025, \textit{MNRAS, 537, 1561}
  \\ Contribution: AGN classification

  \item \textbf{``Misaligned external gas acquisition boosts central black hole activities''} \\
  Yuren Zhou, [4 authors], {\bf Gaoxiang Jin}, et al. 2025, \textit{ApJ, 994, 9}
  \\ Contribution: provide radio AGN classification

  \item \textbf{``Spectroscopically Identified Emission Line Galaxy Pairs in the WISP Survey''} \\
  Y. Sophia Dai, [12 authors], {\bf Gaoxiang Jin} et al. 2021, \textit{ApJ, 923, 156}
  \\ Contribution: figure plotting

  \item \textbf{``Searching for Low-redshift Faint Galaxies with MMT/Hectospec''} \\
  Cheng Cheng, [14 authors], \textbf{Gaoxiang Jin}, et al. 2021, \textit{ApJS, 256, 4}
  \\ Contribution: tests on the source extraction

  \item \textbf{``A Complete 16 $\mu$m Selected Galaxy Sample at z$\sim$1: Mid-infrared Spectral Energy Distributions''} \\
  Jia-Sheng Huang, [21 authors], {\bf Gaoxiang Jin}, et al. 2021, \textit{ApJ, 912, 161}
  \\ Contribution: metallicity estimation and data analysis

  \item \textbf{``A Statistical Study of the Magnetic Imprints of X-class Solar Flares''} \\
  Zekun Lu, Weiguang Cao, \textbf{Gaoxiang Jin}, et al. 2019, \textit{ApJ, 876, 133}
  \\ Contribution: time-series analysis on the magnetic field strength, figure plotting.

  \item \textbf{``What factors shape the radio luminosity of star-forming galaxies? A new calibration from LoTSS-DR2''} \\
  Shravya Shenoy, [8 authors], {\bf Gaoxiang Jin}, et al. \textit{Submitted to MNRAS}
  \\ Contribution: calibrate the single-fiber SFR measurement with the IFU data


  
\end{enumerate}

\vspace{10pt}
\begin{center}
\colorbox{primarycolor!10}{\parbox{0.9\textwidth}{\centering
\color{primarycolor}\faIcon{link} \textbf{Full publication list:} \href{https://ui.adsabs.harvard.edu/user/libraries/SsmUZ6A8RSSVci0wvxpnww}{NASA ADS Library}
}}
\end{center}



\end{document}